\documentclass[11pt]{article}
\usepackage[a4paper,margin=26mm]{geometry}
\usepackage{amsmath,amssymb}
\usepackage[hidelinks]{hyperref}
\setlength{\emergencystretch}{3em}
\title{Conjecture 00000004858: the missing non-elementary hypothesis}
\author{lizaixi01, with OpenAI Codex assistance}
\date{23 October 2026}
\begin{document}
\maketitle
\section{The necessary existence claim}
The source defines a free subgroup as a rank-two nonabelian free subgroup
and proposes a construction in a torsion-free hyperbolic group. It does
not assume that the group is non-elementary. Under this literal universal
construction claim, its necessary existence conclusion is false:
the infinite cyclic group $(\mathbf Z,+)$ is torsion free and hyperbolic,
but has no nonabelian free subgroup.

This refutation concerns that unqualified universal construction. It
does not refute a statement restricted to non-elementary groups, or a
conditional problem in which an existing rank-two free subgroup is
already supplied. The proposed length and algorithmic bounds cannot
repair the missing existence conclusion.

We use the standard word-hyperbolic definition: a finitely generated
group is hyperbolic when its word metric satisfies Gromov hyperbolicity.
For the Cayley-graph formulation and the distinction between elementary
and non-elementary groups, see Dru\c{t}u and Kapovich,
\href{https://www.math.ucdavis.edu/~kapovich/EPR/ggt.pdf}
{\emph{Geometric Group Theory}}, author manuscript, Definition 11.120
and Example 11.121, printed p.407 (PDF p.427). In particular, the latter
example includes free groups of rank one. We prove the required word
metric calculation and hyperbolicity directly.

\section{The genuine shortest-word metric}
Take the finite generating set $\{1\}$ of the additive integer group,
whose signed letters are $1$ and $-1$. A word is an actual finite list
$w=(a_1,\ldots,a_n)$ of these integers, evaluated by its integer sum.
For every such word,
\[
 \left|\sum_{j=1}^n a_j\right|\le n.
\]
Conversely, every integer $z$ is the sum of $|z|$ copies of its sign,
with the empty word for $z=0$. Thus the shortest word length is
\[
 L(z)=\min\{\operatorname{length}(w):
       a_j\in\{1,-1\},\ \operatorname{sum}(w)=z\}=|z|.
\]
Lean defines this length as the actual infimum of the displayed set,
proves nonemptiness, the universal lower bound, and an attaining word.
It then derives the equality, instead of assigning $|z|$ as the
definition of word length. The actual Mathlib Cayley graph for $\{1\}$
has adjacency precisely $y-x=\pm1$.

\newpage
The associated word distance is consequently
\[
 d(x,y)=L(y-x)=|y-x|.
\]
Lean also identifies its real coercion with the actual metric-space
distance on $\mathbf Z$. In particular it is a genuine metric.
For all four integers $x,y,z,t$, direct ordered arithmetic proves
\[
 d(x,y)+d(z,t)\le
 \max\{d(x,z)+d(y,t),\ d(z,y)+d(x,t)\}.
\]
This is Gromov's four-point condition with $\delta=0$.
For the precise metric-space definition, see Nica and \v{S}pakula,
\href{https://arxiv.org/abs/1408.0250v2}{\emph{Strong hyperbolicity}},
arXiv:1408.0250v2, Definition 2.1, printed p.2. We cite only that
metric-space definition: their later standing convention restricting
hyperbolic \emph{groups} to non-elementary ones is explicit on p.3,
and is absent from the competition statement.

The integer group is torsion free: for every nonzero natural number
$n$, the equation $nz=0$ forces $z=0$. The final formal counterexample
collects finite generation, the actual shortest-word hyperbolicity,
torsion freeness, and the subgroup obstruction for the same integer
group. No geometric property is assumed merely by naming a predicate.

\section{No actual free subgroup of rank two}
The standard free group $F_2=\operatorname{FreeGroup}(\mathrm{Fin}\,2)$
has canonical generators $a,b$. Their reduced words show
\[
 ab\ne ba.
\]
For any group homomorphism $f:F_2\to(\mathbf Z,+)$, however, commutativity
of the target gives $f(ab)=f(ba)$. Hence $f$ is not injective.

The Lean target uses \texttt{Multiplicative} to express the actual
additive integer group in multiplicative notation. It separately
proves that, for every actual subgroup $H$ of this group, there is no
group isomorphism $F_2\simeq H$: composing such an isomorphism with the
subgroup inclusion would give the impossible injective homomorphism.
Thus no pair of words, regardless of length or computation time, can
produce the subgroup required by the literal universal claim.

\section{Verification}
The project pins Lean 4.33.0 and Mathlib commit
\begin{center}\small\texttt{db584cd6d46c92f209a44c0f1c829460d327499d}\end{center}
The submission includes a new root build, transitive axiom audits for
every theorem, and full official \texttt{leanchecker --fresh Main}
replay with source-bound receipts. No \texttt{sorry},
\texttt{native\_decide}, or additional axiom is used. OpenAI Codex
assisted, and a separate supervising Codex agent reviewed source,
report and PDF; this is AI review, not human certification.
Organizer acceptance is pending.
\end{document}
